\documentclass[fleqn,10pt]{letter}
\usepackage[utf8]{inputenc}
\usepackage{units}
\usepackage{makecell}
\usepackage{float}
\usepackage[T1]{fontenc}
\usepackage[a4paper, margin=0.87in]{geometry}

\begin{document}

\begin{flushright}
Dr Carmen Cabrera\\
University of Liverpool\\
Liverpool, UK\\
9th of June, 2025\\
\end{flushright}

\begin{flushleft}
To Dr William Burnside \\
Chief Editor \\
\textit{Nature Cities} \\
\end{flushleft} 

Dear Dr Burnside,

I am writing on behalf of my co-authors, as we would greatly appreciate it if you could consider our manuscript titled "Sustained changes to urban mobility after COVID-19 amplified socio-economic inequalities in Latin America" for publication in \textit{Nature Cities}.

Our study presents a large-scale, longitudinal analysis of urban mobility recovery in Latin American countries following the initial COVID-19 disruptions. Using over 170 million anonymised mobile phone location records from Meta-Facebook users (March 2020–May 2022), we investigate how mobility patterns evolved over time, and how these changes intersect with socioeconomic and spatial inequalities.

Prior research has focused on short-term mobility declines in 2020, mostly in high-income countries. Our work fills a critical gap by extending the timeline through 2022, centring the analysis on cities in Latin America. The importance of this work lies on the fact that we find evidence for patterns that have long-term implications for urban resilience and social sustainability, in a region of the world marked by high levels of urbanisation and deep socioeconomic inequalities. Particularly, we find a persistent reduction in net mobility flows in dense urban cores. According to our models, in some cases, these net flows never recover to pre-pandemic baseline levels, suggesting a weakened pull of cities as daily centres of activity. Our findings also indicate that the COVID-19 pandemic has contributed to amplifying pre-existing socioeconomic inequalities, through differences in the recovery of mobility across population groups. 

We consider \textit{Nature Cities} is the ideal venue for our manuscript, due to the close alignment of our research with the journal’s aims and scope. First, we provide new evidence on how pandemic shocks have reshaped urban mobility, hence improving our understanding about the future of cities. Second, our innovative quantitative methodology includes an approach to mitigate representation biases in GPS data, which improves the visibility of marginalised populations in data-sparse contexts. Third, our work contributes directly to the journal’s commitment of featuring more content about the Global South.

%Our findings can help us understand how sustainable our cities are, which is an especially important issue given the rapid urbanisation process that the world is undergoing. Additionally, our results may be used to determine the top-priority regions to be targeted by policies for the alleviation of disruption caused by road accidents. We therefore want to make our work available to a broad, multidisciplinary audience and we believe that publishing in \textit{Nature Cities} can help us achieve this.

We would like to recommend three reviewers for our manuscript based on their expertise in human mobility, urban systems, and data-driven approaches, particularly through the use of location data derived from mobile phone devices:
\begin{itemize}
    \vspace{-1mm}
    \item Professor Esteban Moro; Northeastern University, College of Science
    \vspace{-1mm}
    \item Professor Ronaldo Menezes; University of Exeter, Department of Computer Science
    \vspace{-1mm}
    \item Dr Marta González; University of California, Berkeley, College of Environmental Design
\end{itemize}
\vspace{-1mm}
We thank you very much for the time invested on our behalf and look forward to hearing back from you.

Yours sincerely,

Dr Carmen Cabrera\\
Lecturer in Geographic Data Science\\
University of Liverpool

\end{document}